\textbf{Phần 1. Trắc nghiệm lựa chọn}
\setcounterpageref{ex}{0}

\begin{ex} %[2D13V64]
Để biết số lượng học sinh thích các môn thể thao khác nhau trong lớp, bước đầu tiên em cần làm là
\choice
{Vẽ biểu đồ}
{\True Thu thập số liệu}
{Tính điểm trung bình}
{Xếp hạng học sinh}
\loigiai{
Bước đầu tiên là thu thập thông tin, số liệu từ các bạn.
\textbf{Đáp án của câu hỏi là B}
}
\end{ex}

\begin{ex} %[2D13V64]
Khi quan sát một giỏ hoa quả có cam, táo, chuối, việc xếp cam vào một nhóm, táo vào một nhóm, chuối vào một nhóm gọi là gì?
\choice
{\True Phân loại}
{Kiểm đếm}
{Tính toán}
{Dự đoán}
\loigiai{
Chia các đối tượng theo từng loại đặc điểm gọi là phân loại.
\textbf{Đáp án của câu hỏi là A}
}
\end{ex}

\begin{ex} %[2D13V64]
Trong vườn có $12$ con bướm trắng, $8$ con bướm vàng và $5$ con bướm nâu. Tổng số bướm quan sát được là
\choice
{$24$ con}
{$26$ con}
{\True $25$ con}
{$20$ con}
\loigiai{
$12 + 8 + 5 = 25$ (con bướm).
\textbf{Đáp án của câu hỏi là C}
}
\end{ex}

\begin{ex} %[2D13V64]
Kiểm đếm số đồ dùng học tập của An thấy có: $5$ chiếc bút chì, $3$ cục tẩy, $8$ quyển vở. Đồ dùng nào có số lượng nhiều nhất?
\choice
{Bút chì}
{Cục tẩy}
{Cả ba bằng nhau}
{\True Quyển vở}
\loigiai{
Vì $8 > 5 > 3$ nên vở có số lượng nhiều nhất.
\textbf{Đáp án của câu hỏi là D}
}
\end{ex}

\begin{ex} %[2D13V64]
Trong đợt kiểm tra sức khỏe, tổ Một có $4$ bạn cận thị và $6$ bạn không cận thị. Tổ Một có tất cả bao nhiêu bạn?
\choice
{$8$ bạn}
{\True $10$ bạn}
{$12$ bạn}
{$14$ bạn}
\loigiai{
$4 + 6 = 10$ (bạn).
\textbf{Đáp án của câu hỏi là B}
}
\end{ex}

\begin{ex} %[2D13V64]
Một nhóm bạn đếm số xe đi qua cổng trường trong 5 phút: 15 xe máy, 4 xe đạp, 3 xe ô tô. Số xe máy nhiều hơn tổng số xe đạp và ô tô là
\choice
{\True $8$ xe}
{$7$ xe}
{$9$ xe}
{$10$ xe}
\loigiai{
Tổng xe đạp và ô tô: $4 + 3 = 7$ xe. Xe máy nhiều hơn: $15 - 7 = 8$ xe.
\textbf{Đáp án của câu hỏi là A}
}
\end{ex}

\begin{ex} %[2D13V64]
Phân loại các số sau thành hai nhóm chẵn và lẻ: $12, 17, 24, 35, 40, 51$. Nhóm số chẵn gồm mấy số?
\choice
{$2$ số}
{$4$ số}
{\True $3$ số}
{$5$ số}
\loigiai{
Các số chẵn là 12, 24, 40 (3 số). Các số lẻ là 17, 35, 51 (3 số).
\textbf{Đáp án của câu hỏi là C}
}
\end{ex}

\begin{ex} %[2D13V64]
Khi đếm số lượng bóng bay, bạn Mai dùng các vạch gạch (mỗi cụm 5 vạch biểu thị 5 quả). Mai vẽ được 4 cụm vạch và 2 vạch lẻ. Mai đã đếm được bao nhiêu quả bóng?
\choice
{$20$ quả}
{$24$ quả}
{$18$ quả}
{\True $22$ quả}
\loigiai{
$5 \times 4 + 2 = 22$ (quả bóng).
\textbf{Đáp án của câu hỏi là D}
}
\end{ex}

\begin{ex} %[2D13V64]
Trong lớp học có $18$ bạn nam và $16$ bạn nữ. Sau giờ tan học có $10$ bạn nam và $8$ bạn nữ ra về trước. Trong lớp còn lại bao nhiêu bạn?
\choice
{$14$ bạn}
{\True $16$ bạn}
{$18$ bạn}
{$15$ bạn}
\loigiai{
Tổng sĩ số: $18 + 16 = 34$. Đã về: $10 + 8 = 18$. Còn lại: $34 - 18 = 16$ bạn.
\textbf{Đáp án của câu hỏi là B}
}
\end{ex}

\begin{ex} %[2D13V64]
Một khay đựng các hình hình học: 6 hình tam giác, 8 hình vuông, 10 hình tròn. Khẳng định nào sau đây là SAI?
\choice
{\True Hình tam giác có số lượng nhiều nhất}
{Hình tròn có số lượng nhiều nhất}
{Tổng số hình là 24 hình}
{Hình vuông ít hơn hình tròn 2 hình}
\loigiai{
Hình tròn nhiều nhất (10 hình), tam giác ít nhất (6 hình), nên khẳng định 'tam giác nhiều nhất' là SAI.
\textbf{Đáp án của câu hỏi là A}
}
\end{ex}

\begin{ex} %[2D13V64]
Thu thập số lượng hoa trong lọ: Hoa hồng có $15$ bông, hoa cúc có $20$ bông, hoa ly có $5$ bông. Tỉ lệ số bông hoa cúc so với hoa ly gấp mấy lần?
\choice
{$3$ lần}
{$5$ lần}
{\True $4$ lần}
{$2$ lần}
\loigiai{
$20 : 5 = 4$ lần.
\textbf{Đáp án của câu hỏi là C}
}
\end{ex}

\begin{ex} %[2D13V64]
Kiểm đếm số sách trên 3 ngăn: Ngăn 1 có $35$ cuốn, ngăn 2 có $45$ cuốn, ngăn 3 có $20$ cuốn. Ngăn 2 có nhiều hơn ngăn 3 bao nhiêu cuốn sách?
\choice
{$20$ cuốn}
{$15$ cuốn}
{$10$ cuốn}
{\True $25$ cuốn}
\loigiai{
$45 - 20 = 25$ cuốn.
\textbf{Đáp án của câu hỏi là D}
}
\end{ex}

\textbf{Phần 2. Trắc nghiệm đúng sai}
\setcounterpageref{ex}{0}

\begin{ex} %[2D13V64]
Xét các khẳng định về thu thập và phân loại dữ liệu:
\choiceTF
{\True Phân loại giúp sắp xếp các đồ vật hoặc số liệu có cùng đặc điểm vào chung một nhóm.}
{Kiểm đếm chỉ được thực hiện bằng cách cân các đồ vật.}
{\True Trong một bảng thống kê, tổng số liệu từng phần bằng số liệu toàn thể.}
{\True Có thể phân loại học sinh trong lớp theo nhiều tiêu chí như: giới tính, tháng sinh, sở thích.}
\loigiai{
\begin{itemize}
    \item Ý a) Đúng: Định nghĩa phân loại. Đúng.
    \item Ý b) Sai: Kiểm đếm là đếm số lượng, không phải chỉ cân. Sai.
    \item Ý c) Đúng: Đúng.
    \item Ý d) Đúng: Đúng.
\end{itemize}
\begin{itemize}
    \item a) \textbf{Đúng}.
    \item b) \textbf{Sai}.
    \item c) \textbf{Đúng}.
    \item d) \textbf{Đúng}.
\end{itemize}
}
\end{ex}

\begin{ex} %[2D13V64]
Một hộp bi có 14 viên bi xanh, 16 viên bi đỏ và 10 viên bi vàng:
\choiceTF
{\True Tổng số bi trong hộp là 40 viên.}
{\True Bi đỏ có số lượng nhiều nhất.}
{Số bi xanh ít hơn số bi vàng.}
{\True Số bi đỏ nhiều hơn số bi vàng là 6 viên.}
\loigiai{
\begin{itemize}
    \item Ý a) Đúng: 14 + 16 + 10 = 40 viên. Đúng.
    \item Ý b) Đúng: 16 > 14 > 10. Đúng.
    \item Ý c) Sai: Bi xanh có 14, bi vàng có 10 nên bi xanh nhiều hơn. Sai.
    \item Ý d) Đúng: 16 - 10 = 6 viên. Đúng.
\end{itemize}
\begin{itemize}
    \item a) \textbf{Đúng}.
    \item b) \textbf{Đúng}.
    \item c) \textbf{Sai}.
    \item d) \textbf{Đúng}.
\end{itemize}
}
\end{ex}

\textbf{Phần 3. Trả lời ngắn}
\setcounterpageref{ex}{0}

\begin{ex} %[2D13V64]
Trong vườn có 15 cây cam, 18 cây bưởi và 12 cây chanh. Hỏi trong vườn có tất cả bao nhiêu cây ăn quả?
\shortans{45}
\loigiai{
\textbf{Đáp số:} $45$\\
$15 + 18 + 12 = 45$ cây.
}
\end{ex}

\begin{ex} %[2D13V64]
Một hộp đựng 25 chiếc bút gồm bút bi xanh và bút bi đen, trong đó có 16 bút bi xanh. Hỏi có bao nhiêu bút bi đen?
\shortans{9}
\loigiai{
\textbf{Đáp số:} $9$\\
$25 - 16 = 9$ chiếc.
}
\end{ex}

\begin{ex} %[2D13V64]
Lớp 2A có 36 học sinh. Khi khảo sát môn thể thao yêu thích: 15 bạn thích bóng đá, 12 bạn thích bơi lội, còn lại thích cầu lông. Có bao nhiêu bạn thích cầu lông?
\shortans{9}
\loigiai{
\textbf{Đáp số:} $9$\\
$36 - 15 - 12 = 9$ bạn.
}
\end{ex}

\begin{ex} %[2D13V64]
Kiểm đếm số xe đạp trong bãi: dãy A có 24 xe, dãy B có 28 xe. Cả hai dãy có bao nhiêu xe đạp?
\shortans{52}
\loigiai{
\textbf{Đáp số:} $52$\\
$24 + 28 = 52$ xe.
}
\end{ex}

\begin{ex} %[2D13V64]
Trong đĩa có 8 quả chuối, 6 quả táo và 10 quả dâu tây. Có tất cả bao nhiêu quả?
\shortans{24}
\loigiai{
\textbf{Đáp số:} $24$\\
$8 + 6 + 10 = 24$ quả.
}
\end{ex}

\begin{ex} %[2D13V64]
Một nhóm bạn xếp 5 bàn học, mỗi bàn có 4 bạn ngồi. Đếm được tất cả bao nhiêu bạn?
\shortans{20}
\loigiai{
\textbf{Đáp số:} $20$\\
$4 \times 5 = 20$ bạn.
}
\end{ex}
